\section{Testing and code quality}
\label{sec:testing}

\subsection{Unit tests}

The backend has 148 unit tests, written with \texttt{pytest}, in two families.

\begin{itemize}
  \item \textbf{Service and controller tests} (109 tests) replace the DAOs and the NASA API with
  mocks (\texttt{unittest.mock}). They need neither a database nor the network, run in a few
  seconds, and check the business rules: refusal of invalid values, import split into windows of
  seven days, pause and retry when the NASA rate limit is reached, \emph{partial} import resumed
  the day after, a single import at a time, rights of a user against an administrator\dots
  Controller tests call the endpoints with FastAPI's \texttt{TestClient} and replace the
  logged-in user with \texttt{app.dependency\_overrides}.
  \item \textbf{DAO tests} (39 tests) run on a real PostgreSQL database, in a dedicated schema,
  \texttt{project\_test\_dao}, re-created from \texttt{data/pop\_db\_test.sql} at the start of each
  test session. A guard test checks that they never run on the application schema. Without a
  database, they are skipped instead of failing, so the other tests can still be run anywhere.
\end{itemize}

Most tests follow the \emph{given / when / then} structure and cover both the nominal case and
the error cases (unknown object, duplicate, wrong owner, expired token\dots).

\subsection{Coverage}

The coverage is measured on the service and DAO layers, where the logic is, with branch coverage:
it reaches \textbf{84\,\%} with all the tests. The lines that are not covered are mainly the
\texttt{except} blocks of the DAOs that only log a database error before raising it again, and the
loop of the daily scheduler, which waits 24 hours.

\begin{lstlisting}
cd backend
uv run --all-extras coverage run -m pytest
uv run --all-extras coverage report
\end{lstlisting}

\subsection{Continuous integration and quality tools}

At each push, a GitHub Actions pipeline (\texttt{.github/workflows/ci.yml}) starts a PostgreSQL
service, installs the dependencies, runs all the tests with coverage, then checks the backend and
the frontend with two tools:

\begin{itemize}
  \item \textbf{Ruff}, for the code style (PEP~8, imports, line length of 100) and the formatting;
  formatting on save is enabled in the VSCode workspace;
  \item \textbf{mypy}, which checks that the type annotations are consistent (for instance that a
  value that can be \texttt{None} is not used as if it could not);
  \item \textbf{Bandit}, which looks for common security issues (SQL built by string
  concatenation, hard-coded passwords, requests without timeout\dots).
\end{itemize}

Every class and public method has a docstring, and the logs record each HTTP request with its
status and the inputs and outputs of the main methods, with the passwords masked.

\subsection{Final review}

Before the code freeze, the whole application was run end to end on a fresh database and the code
was reviewed. This review found several bugs, which were fixed, each with a test:

\begin{itemize}
  \item an expired access token was still accepted, because the token was looked up in the
  connection history after its signature check had failed; the token could also be passed in the
  URL, where it ends up in the logs;
  \item the search and the NEO details could be read without being logged in;
  \item adding a non-existent NEO to the favourites crashed with an error 500;
  \item a NEO could be declared with negative diameters or a minimum above the maximum, and a
  validation error of this kind crashed the error handler itself;
  \item NEOs declared without a diameter were counted among the objects larger than 1~km;
  \item the same notification was shown at every login;
  \item the frontend never used the refresh token, so users were logged out after 30 minutes, and
  the daily update was only triggered by the frontend, under a condition that blocked it for a
  year after the first import.
\end{itemize}
